\section{Remote sequences and local persistence}
\label{sec:arch-remote-sequences}

An installation built with external-provider support exposes metadata,
complete nucleotide records and sequence chunks through relations in the
\texttt{bioseq\_external} schema. These relations obtain their rows from remote
services. An application can select a versioned accession and insert its
sequence into a local BioSeqDB table before repeated analysis.

Metadata discovery and sequence retrieval use distinct requests. Generated
HTTPS clients perform bounded metadata requests, and provider adapters
normalize their responses. For NCBI, the catalog reads Datasets assembly
reports~\autocite{ncbi_datasets_v2_api}. A nucleotide accession such as
\texttt{NC\_012920.1} instead requests plain-text FASTA from the EFetch
nucleotide endpoint~\autocite{ncbiEutilities}.
Assembly accessions beginning with \texttt{GCF\_} or \texttt{GCA\_} select a
Datasets ZIP package containing sequence files. An assembly accession and a
constituent sequence accession therefore identify different retrieval units.

The remote relation returns an \emph{unbound sequence literal}: sequence
characters without a selected storage alphabet or packed layout. In
Listing~\ref{lst:arch-ncbi-import}, assignment to the destination
\texttt{bioseq.seq('dna5',64)} column validates and encodes those characters
for the declared alphabet and storage width.

\begin{figure}[htbp]
  \centering
  \begin{tikzpicture}[
    x=1cm,y=1cm,>=Latex,font=\footnotesize,
    archremote/.style={draw=black!55,rounded corners=2pt,
      fill=blue!5,align=center,inner sep=4pt,minimum height=.70cm},
    archflow/.style={->,semithick,draw=black!75}
  ]
    \node[font=\footnotesize\bfseries] at (0,1.2) {Remote service};
    \node[font=\footnotesize\bfseries] at (4.70,1.2) {Network worker};
    \node[font=\footnotesize\bfseries] at (9.65,1.2) {PostgreSQL backend};
    \draw[densely dashed,black!35] (1.60,.62) -- (1.60,-.68);
    \draw[densely dashed,black!35] (8.05,.62) -- (8.05,-2.4);
    \node[archremote,text width=2.10cm] (archremote-ncbi) at (0,0)
      {NCBI EFetch\\FASTA response};
    \node[archremote,text width=2.15cm] (archremote-reader) at (3.25,0)
      {Read and parse\\sequence bases};
    \node[archremote,text width=2.05cm] (archremote-queue) at (6.40,0)
      {Bounded queue\\owned sequence data};
    \node[archremote,text width=2.60cm] (archremote-backend) at (9.65,0)
      {Build unbound SQL value\\apply local checks};
    \node[archremote,text width=2.60cm,fill=green!7]
      (archremote-local) at (9.65,-1.9)
      {Local table\\packed sequence};
    \draw[archflow] (archremote-ncbi) -- (archremote-reader);
    \draw[archflow] (archremote-reader) -- (archremote-queue);
    \draw[archflow] (archremote-queue) -- (archremote-backend);
    \draw[archflow] (archremote-backend) --
      node[left,font=\footnotesize] {\texttt{INSERT} + destination type}
      (archremote-local);
  \end{tikzpicture}
  \caption{An exact NCBI sequence request followed by local insertion.
  Remote delivery and persistent sequence encoding have separate owners.}
  \label{fig:arch-remote-import}
\end{figure}

\noindent\begin{minipage}{\linewidth}
\begin{lstlisting}[style=bioseq-arch-sql,basicstyle=\footnotesize\ttfamily,
  caption={Importing one versioned NCBI sequence into local storage.},
  label={lst:arch-ncbi-import}]
CREATE TABLE ncbi_sequences (
    accession text,
    sequence bioseq.seq('dna5', 64),
    retrieved_at timestamptz
) USING bioseq;

INSERT INTO ncbi_sequences
SELECT accession, sequence, retrieved_at
FROM bioseq_external.sequences
WHERE provider = 'ncbi' AND accession = 'NC_012920.1';

SELECT accession, bioseq.len(sequence) AS bases
FROM ncbi_sequences;
\end{lstlisting}
\end{minipage}

Ordinary SQL qualifications remain local checks even when equality predicates
also constrain the remote request. Array expansion is capped at 64 query
combinations, with a warning when combinations are omitted. A provider union
may return partial results if one provider fails, and metadata scans do not
automatically follow every page. Explicit accessions and a single provider
make the import's requested scope clear. Selecting fewer columns from a
sequence relation does not currently avoid sequence retrieval.

For the NCBI streaming route, a worker passes owned sequence data through a queue
bounded by both bytes and item count. Chunk delivery limits temporary
buffering, whereas a complete-record value still requires record-sized memory.
Archive retrieval buffers a size-limited ZIP before decoding. The backend checks
cancellation between bounded queue waits. Transient FASTA requests can be
retried only before the first row is emitted, preventing replay of visible
chunks.

Remote responses are not a durable cache, and request rates are not coordinated
across backends. E-utilities usage requirements therefore remain relevant to
the caller~\autocite{ncbiEutilitiesUsage}. A transaction controls the local INSERT,
but does not make separate HTTP requests one database snapshot. Retaining the
accession version and retrieval time records which import supplied the stored
sequence. Subsequent queries of the local table need no remote request.
